\documentclass[../../../include/open-logic-section]{subfiles}

\begin{document}

\olfileid{sfr}{infinite}{card-sb}

\olsection{附录：证明 Schr\"oder-Bernstein 定理}

在我们告别朴素集合论之前，还要考虑一个朴素（却精巧的！）证明。这就是 Schr\"oder-Bernstein 定理（\olref[sfr][siz][sb]{thm:schroder-bernstein}）的证明：若 $\cardle{A}{B}$ 且 $\cardle{B}{A}$，则 $\cardeq{A}{B}$；换言之，给定单射 $f \colon A \to B$ 和 $g \colon B \to A$，存在一个双射 $h \colon A \to B$。

在本章中，我们遵循了 Dedekind（戴德金）的\emph{闭包}概念。事实上，Dedekind 利用这个概念给出了 Schr\"oder-Bernstein 定理的一个优美证明，我们将在此介绍。该证明紧随 \citet[pp.~157--8]{Potter2004} 的论述；若你想要略有不同但本质相似的处理，可以参考该书。稍作搜索你也会发现，这个定理——很像 $\sqrt{2}$ 的无理性——有许多有趣且不同的证明。

使用与 \olref[sfr][infinite][dedekind]{Closure} 类似的记号，对每个集合 $B$ 和函数~$f$，令
\[
\Closureofunder{f}{B} = \bigcap \Setabs{X}{B \subseteq X 
\text{ 且 $X$ 是 $f$-闭的}}
\]。
这样定义的 $\Closureofunder{f}{B}$ 是包含~$B$ 的最小 $f$-闭集，即：

\begin{lem}\ollabel{Closureprops}
	对任意函数 $f$ 和任意 $B$：
	\begin{enumerate}
		\item\ollabel{Closurehaselem} $B \subseteq \Closureofunder{f}{B}$；且
		\item\ollabel{Closureclosed} $\Closureofunder{f}{B}$ 是 $f$-闭的；且
		\item\ollabel{Closuresmallest} 若 $X$ 是 $f$-闭的且 $B \subseteq X$，则 $\Closureofunder{f}{B} \subseteq X$。
	\end{enumerate}
\end{lem}

\begin{proof}
	与 \olref[sfr][infinite][dedekind]{closureproperties} 完全相同。
\end{proof}


我们需要最后一个事实来得出 Schr\"oder-Bernstein 定理：

\begin{prop}\ollabel{sbhelper}
若 $A \subseteq B \subseteq C$ 且 $\cardeq{A}{C}$，则 $\cardeq{A}{B}$ 且~$\cardeq{B}{C}$。
\end{prop}

\begin{proof}
	给定一个双射 $f \colon C \to A$，令 $F =
\Closureofunder{f}{C \setminus B}$ 并按如下方式定义函数 $g$，其定义域为 $C$：
	\[
	g(x) = 
	\begin{cases}
			f(x) &\text{若 $x \in F$}\\
			x & \text{否则}
		\end{cases}
\]
	下面证明 $g$ 是从 $C$ 到 $B$ 的双射（即 $\cardeq{B}{C}$）。由此可知 $\comp{f^{-1}}{g}
\colon A \to B$ 是双射，从而完成证明（$\cardeq{A}{B}$）。

	首先断言：若 $x \in F$ 但 $y \notin F$，则 $g(x) \neq g(y)$。反设不然，即 $y = g(y) = g(x) =
f(x)$。由于 $x \in F$，且由 \olref{Closureprops} 知 $F$ 是 $f$-闭的，故 $y = f(x) \in F$，矛盾。

	现假设 $g(x) = g(y)$。那么，由上述断言，$x \in F$ 当且仅当 $y \in F$。
	如果 $x, y \in F$，则 $f(x) = g (x) = g(y) = f(y)$，由于 $f$ 是双射，可得 $x = y$。
	如果 $x, y \notin F$，则 $x = g(x) =
g(y) = y$。所以 $g$ 是单射。

	还须证明 $\ran{g} = B$。取定 $x \in B \subseteq C$。
	若 $x \notin F$，则 $g(x) = x$。若 $x \in F$，则存在某个 $y \in F$ 使得 $x = f(y)$，否则 $F \setminus \{x\}$ 将是 $f$-闭的且包含 $C\setminus B$，由 \olref{Closureprops} 可知这是不可能的；进而 $g(y) = f(y) = x$。
\end{proof}


最后，以下是主要结果的证明。回顾一下：给定函数 $h$ 和集合 $D$，我们定义 $\funimage{h}{D} = \Setabs{h(x)}{x
\in D}$。

\begin{proof}[Schr\"oder-Bernstein 定理的证明]
	设 $f \colon A \to B$ 和 $g \colon B \to A$ 为单射。由于 $\funimage{f}{A} \subseteq B$，我们有 $\funimage{g}{\funimage{f}{A}} \subseteq
g[B] \subseteq A$。此外，$\comp{f}{g} \colon A \to
\funimage{g}{\funimage{f}{A}}$ 是单射，因为 $g$ 和 $f$ 都是单射；而且 $\comp{f}{g}$ 实际上是双射，这由我们对其陪域的定义即可看出。所以
	$\cardeq{\funimage{g}{\funimage{f}{A}}}{A}$，从而根据 \olref{sbhelper}，存在一个双射 $h \colon A \to
\funimage{g}{B}$。再者，$g^{-1}$ 是从 $\funimage{g}{B}$ 到 $B$ 的双射。因此 $\comp{h}{g^{-1}} \colon A \to B$ 是双射。
\end{proof}

\end{document}